%% kapa: exam=1 round=35 year=2024 period=1 form=A subject=economics subject_ko=경제학원론 seq=38 number=78
%% source: Q-Net 감정평가사 기출문제 (cst00309) artlSeq=5213229 file=raw/1008/5213229/1교시 문제지 원본.pdf
%% conversion: kordoc 4.24.51 -> tools/gen1.py (auto); review=proofread-sonnet; flags=none
\begin{question}{78}
거시경제이론과 관련된 경제학파에 대한 설명으로 옳은 것은?
\begin{choices}
  \item 새케인즈학파(New Keynesian)는 단기 필립스곡선이 수직이라고 주장한다.
  \item 새케인즈학파는 가격 신축성에 근거하여 경기변동을 설명한다.
  \item 새케인즈학파는 단기에서 화폐중립성이 성립한다고 주장한다.
  \item 실물경기변동이론에 따르면 경기변동국면에서 소비의 최적화가 달성된다.
  \item 새고전학파는 메뉴비용의 존재가 경기변동에 중요한 역할을 한다고 주장한다.
\end{choices}
\end{question}
\answer{4}
